% TOP_PROBLEM: {"id": 30003585, "title": "Toms–Winter Regularity Conjecture", "category_id": 4, "category": {"id": 4, "name": "algebra", "display_name": "Algebra", "description": "Group theory, ring theory, field theory, and algebraic structures.", "slug": "algebra", "order_index": 4, "created_at": "2026-07-31T15:26:25.671Z"}, "status": "partially_solved"}
% FIRST_POSTED: 2026-10-06
% ATTEMPT_STATUS: unresolved
% OPENAI_CLAIM: {"schema_version":1,"source_repo":"https://github.com/openai/math","source_commit":"adc7f1241b42e322a6451854ab7e4b4c146bf78a","release_date":"2026-10-06","match":"partial","resolution":"partial","independently_reviewed":false,"second_pass":"n/a","families":[{"family":"291","title":"Cuntz comparison, nuclear dimension, and equivariant Jiang–Su stability","subject":"Operator algebras","manuscripts":[{"title":"Cuntz comparison and Jiang–Su absorption","pdf_path":"preprints/Cuntz-comparison-and-Jiang-Su-absorption-September-23-2026/paper.pdf","readme_path":"preprints/Cuntz-comparison-and-Jiang-Su-absorption-September-23-2026/README.md","date":"2026-09-23","theorem_ref":"Introduction main regularity theorems"},{"title":"Nuclear dimension and Jiang–Su stability without elementary subquotients","pdf_path":"preprints/Nuclear-dimension-and-Jiang-Su-stability-without-elementary-subquotients-September-23-2026/paper.pdf","readme_path":"preprints/Nuclear-dimension-and-Jiang-Su-stability-without-elementary-subquotients-September-23-2026/README.md","date":"2026-09-23","theorem_ref":"Introduction main regularity theorems"},{"title":"Tracial projection methods and uniform property Gamma","pdf_path":"preprints/Tracial-projection-methods-and-uniform-property-Gamma-September-23-2026/paper.pdf","readme_path":"preprints/Tracial-projection-methods-and-uniform-property-Gamma-September-23-2026/README.md","date":"2026-09-23","theorem_ref":"Introduction main regularity theorems"}],"lean":{"doc_path":"lean/docs/291.md","comparators":["lean/ComparatorChallenges/UniformGamma.json","lean/ComparatorChallenges/UniformGamma.lean"],"declarations":["OAI.ComparatorModel.CurrentMain.main"],"covers_main_theorem":false},"reasoning_trace":null},{"family":"302","title":"Radius of comparison equals half the mean dimension","subject":"Operator algebras","manuscripts":[{"title":"Radius of comparison equals half the mean dimension","pdf_path":"preprints/Radius-of-comparison-equals-half-the-mean-dimension-September-25-2026/paper.pdf","readme_path":"preprints/Radius-of-comparison-equals-half-the-mean-dimension-September-25-2026/README.md","date":"2026-09-25","theorem_ref":"preprints/Radius-of-comparison-equals-half-the-mean-dimension-September-25-2026/build/sections/01-introduction.tex#thm:main"}],"lean":null,"reasoning_trace":null}]}
\documentclass[11pt]{article}
\usepackage[margin=25mm]{geometry}
\usepackage{fontspec}
\setmainfont{Latin Modern Roman}
\usepackage{amsmath,amssymb,mathtools}
\usepackage{unicode-math}
\setmathfont{Latin Modern Math}
\usepackage{enumitem}
\usepackage{xurl,hyperref,bookmark}
\hypersetup{colorlinks=true,urlcolor=blue}
\setcounter{secnumdepth}{0}
\setlength{\parindent}{0pt}
\setlength{\parskip}{5pt}
\setlength{\emergencystretch}{3em}
\begin{document}
\section*{Toms–Winter Regularity Conjecture}

\subsection{OpenAI's claim}
Manuscript 1 (family 291, 2026-09-23): A separable nuclear C*-algebra that is simple and non-elementary is claimed to absorb Jiang--Su under extended-functional strict comparison, including nonunital and unbounded-trace cases; semigroup divisibility also yields a nonsimple criterion; Manuscript 2 (family 291, 2026-09-23): Finite nuclear dimension is claimed to characterize Jiang--Su stability for separable nuclear algebras lacking elementary ideal subquotients; moreover, Jiang--Su stabilization always lowers nuclear dimension to at most one. Manuscript 3 (family 291, 2026-09-23): The paper claims uniform Gamma from real rank zero in the tracial-completion ultrapower for simple unital nuclear stably finite algebras. Comparison gives absorption, also under specified bounded-trace hypotheses in the stably projectionless setting; Manuscript 4 (family 302, 2026-09-25): Whenever an infinite compact metrizable system is minimal under a homeomorphism, its crossed-product comparison radius is claimed to be half its mean dimension, including infinite values; zero forces nuclear dimension at most one.

\subsection{Outline of the argument}
Manuscript 1 (family 291, 2026-09-23): Trace transfer into small hereditary supports followed by spectral rounding establishes full Cuntz-semigroup divisibility. Pure-state compression produces norm-visible matrix cones, while localized nuclear approximations produce trace-visible ones. Comparison packs their defects into a central corner, constructing dimension-drop maps whose lifted intertwining relations give Jiang--Su absorption. Manuscript 2 (family 291, 2026-09-23): Localized order-zero approximations control ranks even when traces are infinite on some ideals. A supporting map and comparison yield two-color approximations that survive stabilization and hereditary passage. Conversely, Gaussian sign tensors and pure-state excision construct two full orthogonal central elements from the absence of characters, supplying the absorption criterion.

Manuscript 3 (family 291, 2026-09-23): Nuclear averaging and scalar block compressions preserve coefficient traces while packing small projections with controlled commutators. A separate comparison construction rounds ranks, forms clipped sums of centered symmetries, and extracts Gaussian trace partitions. In the projectionless case, a combined overlap-and-model cost controls small splittings before reindexing amplifies them to exact central partitions. Manuscript 4 (family 302, 2026-09-25): Finite-support comparison witnesses are converted to bundle embeddings, which feed the companion's boundary-preserving compression theorem. Packet reconstruction translates that compression into a smaller cover without introducing absent colors. Applying the construction on aligned orbit blocks contradicts any larger cover-growth rate, providing the missing lower bound; established upper bounds and regularity results finish the equality and zero-dimensional consequences.

\subsection{Scope relative to this problem}
Manuscript 1 (family 291, 2026-09-23): The strict-comparison/Jiang–Su/nuclear-dimension claims do not settle the additional finite-decomposition-rank equivalence for all stably finite algebras requested by this record.

Manuscript 2 (family 291, 2026-09-23): The strict-comparison/Jiang–Su/nuclear-dimension claims do not settle the additional finite-decomposition-rank equivalence for all stably finite algebras requested by this record.

Manuscript 3 (family 291, 2026-09-23): The strict-comparison/Jiang–Su/nuclear-dimension claims do not settle the additional finite-decomposition-rank equivalence for all stably finite algebras requested by this record.

Manuscript 4 (family 302, 2026-09-25): The primary manuscript establishes radius of comparison equals half mean dimension for every minimal infinite compact-metrizable Z system, with zero-dimensionality regularity equivalences. The result proves the radius/mean-dimension identity and comparison/absorption/nuclear-dimension relations for minimal Z crossed products C(X) crossed\_h Z. It does not quantify over every simple nuclear algebra, nonunital algebra, or all decomposition-rank clauses of the selected statement.

\subsection{Formal verification}
For family 291, the scope page documents the following selections. The selection gives uniform Gamma from real rank zero of the designated tracial ultrapower. The documented comparison scope also treats Jiang--Su absorption and Gamma under strict comparison, with additional trace and target-rank conditions in the projectionless case. Manuscript 1 is not among the papers selected by that page, so the family link does not formalize its main conclusion. It does not establish formal coverage of the complete problem clause. Manuscript 2 is not among the papers selected by that page, so the family link does not formalize its main conclusion. It does not establish formal coverage of the complete problem clause. The documented statement covers the main claim of manuscript 3. It does not establish formal coverage of the complete problem clause. Comparator files: UniformGamma.lean. Declarations: OAI.ComparatorModel.CurrentMain.main.

Family 302 has no Lean scope page in this snapshot; no formalization of these manuscript results is listed.

This site has not executed the Lean files or checked their dependencies. These descriptions report the scope of the linked declarations, rather than a new verification of the proofs.

\subsection{Status caveat}
OpenAI attributes these results to an internal model. Its release notice cautions that results lacking full formalization can contain mistakes. This site has not independently assessed the proofs; a solved label records the scope of the reported claim.

\subsection{Sources}
\begin{itemize}
\item OpenAI, \emph{Cuntz comparison and Jiang–Su absorption}, 2026-09-23; citation key \texttt{OAI:Cuntz-comparison-and-Jiang-Su-absorption-September-23-2026}. \href{https://github.com/openai/math/blob/adc7f1241b42e322a6451854ab7e4b4c146bf78a/preprints/Cuntz-comparison-and-Jiang-Su-absorption-September-23-2026/paper.pdf}{PDF (GitHub)}; \href{https://github.com/openai/math/raw/adc7f1241b42e322a6451854ab7e4b4c146bf78a/preprints/Cuntz-comparison-and-Jiang-Su-absorption-September-23-2026/paper.pdf}{PDF (direct)}; \href{https://github.com/openai/math/blob/adc7f1241b42e322a6451854ab7e4b4c146bf78a/preprints/Cuntz-comparison-and-Jiang-Su-absorption-September-23-2026/README.md}{Manuscript page and citation}; \href{https://github.com/openai/math/blob/main/preprints/Cuntz-comparison-and-Jiang-Su-absorption-September-23-2026/paper.pdf}{Latest version}.
\item OpenAI, \emph{Nuclear dimension and Jiang–Su stability without elementary subquotients}, 2026-09-23; citation key \texttt{OAI:Nuclear-dimension-and-Jiang-Su-stability-without-elementary-subquotients-September-23-2026}. \href{https://github.com/openai/math/blob/adc7f1241b42e322a6451854ab7e4b4c146bf78a/preprints/Nuclear-dimension-and-Jiang-Su-stability-without-elementary-subquotients-September-23-2026/paper.pdf}{PDF (GitHub)}; \href{https://github.com/openai/math/raw/adc7f1241b42e322a6451854ab7e4b4c146bf78a/preprints/Nuclear-dimension-and-Jiang-Su-stability-without-elementary-subquotients-September-23-2026/paper.pdf}{PDF (direct)}; \href{https://github.com/openai/math/blob/adc7f1241b42e322a6451854ab7e4b4c146bf78a/preprints/Nuclear-dimension-and-Jiang-Su-stability-without-elementary-subquotients-September-23-2026/README.md}{Manuscript page and citation}; \href{https://github.com/openai/math/blob/main/preprints/Nuclear-dimension-and-Jiang-Su-stability-without-elementary-subquotients-September-23-2026/paper.pdf}{Latest version}.
\item OpenAI, \emph{Tracial projection methods and uniform property Gamma}, 2026-09-23; citation key \texttt{OAI:Tracial-projection-methods-and-uniform-property-Gamma-September-23-2026}. \href{https://github.com/openai/math/blob/adc7f1241b42e322a6451854ab7e4b4c146bf78a/preprints/Tracial-projection-methods-and-uniform-property-Gamma-September-23-2026/paper.pdf}{PDF (GitHub)}; \href{https://github.com/openai/math/raw/adc7f1241b42e322a6451854ab7e4b4c146bf78a/preprints/Tracial-projection-methods-and-uniform-property-Gamma-September-23-2026/paper.pdf}{PDF (direct)}; \href{https://github.com/openai/math/blob/adc7f1241b42e322a6451854ab7e4b4c146bf78a/preprints/Tracial-projection-methods-and-uniform-property-Gamma-September-23-2026/README.md}{Manuscript page and citation}; \href{https://github.com/openai/math/blob/main/preprints/Tracial-projection-methods-and-uniform-property-Gamma-September-23-2026/paper.pdf}{Latest version}.
\item Family 291: \href{https://github.com/openai/math/blob/adc7f1241b42e322a6451854ab7e4b4c146bf78a/lean/docs/291.md}{Lean scope}.
\item \href{https://github.com/openai/math/blob/adc7f1241b42e322a6451854ab7e4b4c146bf78a/lean/ComparatorChallenges/UniformGamma.json}{Comparator UniformGamma.json}.
\item \href{https://github.com/openai/math/blob/adc7f1241b42e322a6451854ab7e4b4c146bf78a/lean/ComparatorChallenges/UniformGamma.lean}{Comparator UniformGamma.lean}.
\item OpenAI, \emph{Radius of comparison equals half the mean dimension}, 2026-09-25; citation key \texttt{OAI:Radius-of-comparison-equals-half-the-mean-dimension-September-25-2026}. \href{https://github.com/openai/math/blob/adc7f1241b42e322a6451854ab7e4b4c146bf78a/preprints/Radius-of-comparison-equals-half-the-mean-dimension-September-25-2026/paper.pdf}{PDF (GitHub)}; \href{https://github.com/openai/math/raw/adc7f1241b42e322a6451854ab7e4b4c146bf78a/preprints/Radius-of-comparison-equals-half-the-mean-dimension-September-25-2026/paper.pdf}{PDF (direct)}; \href{https://github.com/openai/math/blob/adc7f1241b42e322a6451854ab7e4b4c146bf78a/preprints/Radius-of-comparison-equals-half-the-mean-dimension-September-25-2026/README.md}{Manuscript page and citation}; \href{https://github.com/openai/math/blob/main/preprints/Radius-of-comparison-equals-half-the-mean-dimension-September-25-2026/paper.pdf}{Latest version}.
\item \href{https://github.com/openai/math/blob/adc7f1241b42e322a6451854ab7e4b4c146bf78a/CONTENTS.md}{OpenAI catalogue entry}.
\end{itemize}
\end{document}
